\section{When did discrete representation get so bad?}

Flow-based policies have emerged as the de facto parameterization of choice for
robot learning, while discrete autoregressive (AR) alternatives are often
dismissed on the grounds that quantizing continuous actions necessarily
sacrifices fidelity. Yet, in other modalities, such as audio, autoregressive
models have proven remarkably effective despite relying on discrete
representations. This contrast raises a natural question: \emph{Is quantization
really what holds back autoregressive policies in robotics?} The answer is
perhaps surprising: quantization error alone is often insufficient to explain
the failure of discrete policies. In the majority of cases, policies fail by a
wide margin and cannot be explained by quantization error alone. Moreover,
discrete actions are in fact more precise in tail-end situations. So why do
discrete policies fail in the first place? As a starting point, we argue that
the primary limitations of autoregressive action generation are uneven
information allocation and error accumulation, rather than quantization error.
In commonly used action representations such as FAST, uneven information
allocation leads to variation in conditional token entropy across positions,
which may result in unstable training dynamics. On the other hand, for action
binning, long autoregressive sequences lead to error accumulation and degraded
performance. To address both problems, we advocate a simple learned tokenization
that encourages even information allocation while maintaining compression.
